\documentclass[10pt,a4paper]{amsart}
\usepackage[margin=19mm]{geometry}
\usepackage{amsmath,amssymb,mathtools}
\usepackage[scr=rsfs,cal=euler]{mathalfa}
\usepackage{xfrac}
\usepackage[unicode,hidelinks,bookmarks=false]{hyperref}
\newcommand{\ie}{i.e.\ }
\newcommand{\cf}{cf.\ }
\newcommand{\resp}{respectively\ }
\newcommand{\Subsection}{\subsection}
\providecommand{\nobreakdash}{\nobreak\hbox{-}}
\setlength{\emergencystretch}{2em}
\multlinegap0pt
\usepackage{bm}
\usepackage{bbm}
\usepackage{dsfont}
\usepackage{xspace}
\usepackage{cancel}
\usepackage{mathtools}
\usepackage[shortlabels]{enumitem}
\usepackage{amsthm}
\makeatletter
\@ifundefined{theorem}{\newtheorem{theorem}{Theorem}}{}
\makeatother
\newcommand{\N}{\mathbb{N}}  % set of natural numbers
\newcommand{\Z}{\mathbb{Z}}  % set of integer numbers
\newcommand{\cF}{\mathcal{F}}
\newcommand{\cantor}{\mathbf C}
\newcommand{\fence}{{\bf F}_{\Phi}{}}
\newcommand{\fstructure}{\{(V_n, \Psi_n,\varphi^L_n,\varphi^U_n)\}_{n=0}^{\infty}}
\newcommand{\fsystem}{\{(G_n, \Psi_n,\varphi^L_n,\varphi^U_n)\}_{n=0}^{\infty}}
\newcommand{\groupcantor}{{\mathbb H}(\cantor)}
\newcommand{\groupfence}{{\mathbb{H}(\fence)}}
\newcommand{\homeocantor}{{\mathbf H}_{X}}
\newcommand{\oa}{\overrightarrow}
\newcommand{\orb}{{\mathrm{Orb}}}
\title[Dynamics on fences]{Dynamics on fences}
\author{Jernej \v{C}in\v{c}, Udayan B. Darji, Benjamin Vejnar}
\date{Source: arXiv:2604.07088v1 (CC BY 4.0)}
\begin{document}
\maketitle

\noindent\emph{Source and licence.} Dynamics on fences. Version arXiv:2604.07088v1 (CC BY 4.0), distributed under the Creative Commons Attribution 4.0 International licence, as recorded in the arXiv licence metadata for this exact version and shown on the abstract page https://arxiv.org/abs/2604.07088v1. Full rights notice: licenses/arxiv-2604.07088-CC-BY-4.0.txt.\par\smallskip

\noindent\textit{Editorial scope.} Counted unit: the Theorem whose source label is \texttt{thm:mixing} in the article (source0.tex lines 1183--1192, source label \texttt{thm:mixing}), delivered together with the article's own complete original proof of it (source0.tex lines 1193--1272, one \texttt{proof} environment of 80 lines). No mathematics is written, repaired or re-derived: statement and proof are the article's own text line for line. Required context retained uncounted: thm:fenceconst (lines 587--604); eq:s(x) (lines 833--836); thm:mainF (lines 861--872); thm:transitive (lines 1055--1062). Every definition, display and label of those lines is retained unchanged. Omitted from this package and not redistributed: 1--586, 605--832, 837--860, 873--1054, 1063--1182, 1273--1465. Nothing inside a retained excerpt is omitted or altered: every retained body is delivered exactly as the article's lines are written, so no omission receipt of any kind accompanies this package. Citations: the retained excerpts carry no citation command at all, so no bibliography entry of the article is transcribed and no reference is redistributed. Reference adaptation: none was needed and none was made; every reference inside the delivered unit resolves inside it, so no unresolved reference or citation is produced. Printed slips kept literal: no printed word, symbol, hypothesis or number of the counted statement or of its proof was altered. Layout and class: the article's own class is \texttt{amsart} with packages including graphicx, inputenc, comment, babel, enumitem, amsmath, amsfonts, amssymb, amsthm, amscd, tikz, geometry, epsfig, color; the wrapper is the lane's pinned \texttt{amsart} editorial wrapper at 10pt on A4, which restates the theorem-like environments the retained text needs and adds the article's own macro definitions that the retained text needs (\texttt{\textbackslash N}, \texttt{\textbackslash Z}, \texttt{\textbackslash cF}, \texttt{\textbackslash cantor}, \texttt{\textbackslash fence}, \texttt{\textbackslash fstructure}, \texttt{\textbackslash fsystem}, \texttt{\textbackslash groupcantor}, \texttt{\textbackslash groupfence}, \texttt{\textbackslash homeocantor}, \texttt{\textbackslash oa}, \texttt{\textbackslash orb}), with extra wrapper packages (bm, bbm, dsfont, xspace, cancel, mathtools, enumitem). The delivered numbering is the unit's own. Assets: no retained span contains a figure environment, a graphics-inclusion command, a PDF-page inclusion, a table or a TikZ picture; no image file is copied into this package, the package declares no asset dependency and no image file is redistributed with it. Licence: version arXiv:2604.07088v1 of this article is distributed under the Creative Commons Attribution 4.0 International licence, as recorded in the arXiv licence metadata for this exact version and shown on the abstract page https://arxiv.org/abs/2604.07088v1. A full rights notice is delivered at licenses/arxiv-2604.07088-CC-BY-4.0.txt. Characters: every retained excerpt is pure ASCII, so the delivered engine, XeLaTeX with its default Latin Modern text font, reports no missing character.
\begin{theorem}\label{thm:fenceconst} Let $\fstructure$ be an $\cF$-structure and $\fence$ be its associated fence. Furthermore, assume that the following condition is satisfied.  

\begin{equation}\tag{$\dagger$}\label{eq:dagger}
\forall v \in V_n,\  \exists v' \in V_{n+1} \text{ such that } \Psi_n(v') = v , \  \varphi^L_n(v) = \varphi^L_n(v') \text{ and } \varphi^U_n(v) = \varphi^U_n(v').\
\end{equation}

    \begin{enumerate}
    \item If $\varphi^L$ is the zero function and $\varphi^U$ is constant one function, then $\fence$ is the Cantor fence.
    \item  If $\{\eta^+_n\}  \rightarrow 0$, then $\fence$ is the Scissorhand Fence.
    \item If $\varphi^L$ is the zero function and $\{\eta^+_n\} \rightarrow 0$, then $\fence$ is the Lelek fence.
    \item If $\{\eta^+_n\} \rightarrow 0$, and $\{\eta^-_n\} \rightarrow 0$ then $\fence$ is a two-sided Scissorhand Fence.
        \item If $\{\eta_n\} \rightarrow 0$, then $\fence$ is the Fra\"iss\'e fence.

        %\item If $\{\eta^+_{2n+1}\} \rightarrow 0$, and $\{\eta^-_{2n+2}\} \rightarrow 0$ then $\fence$ is the Fra\"iss\'e fence.
    
    \end{enumerate}
    
\end{theorem}

\begin{equation}\label{eq:s(x)}
%s(x):=\lim_{n\to \infty} s_n(u_n,v_n) \quad 
s(x):= \lim_{n\to \infty} \tilde{s}_n (x).
\end{equation}

\begin{theorem}\label{thm:mainF}
Let $\fsystem$ be an $\cF$-system which satisfies Condition $\Gamma$ and  ${\bf F}_{\Phi}$ is a Scissorhand Fence. 
Then, there exists a unique continuous surjection $T: {\bf F}_{\Phi}\rightarrow {\bf F}_{\Phi}$, with $\homeocantor$ as a factor, satisfying \[T\left(x, \varphi^U(x)\right)=\left(\homeocantor(x), \varphi^U(\homeocantor(x))\right)\] for all $x \in X$.
 Moreover, 
 \begin{enumerate}
         \item if $\homeocantor$ is a homeomorphism, then $T$ is a homeomorphism of ${\bf F}_{\Phi}$.
     %\item if $\homeocantor$ is minimal and ${\bf F}_{\Phi}$ is a two-sided Scissorhand Fence, then $T$ is minimal.
     \item if $s$ as defined in \eqref{eq:s(x)} and $\homeocantor$ are both Lipschitz, then so is $T$.
     \item if $s$ and $1/s$ are both Lipschitz and $\homeocantor$ is bi-Lipschitz, then $T$ is bi-Lipschitz.
     \item if $s=1$ and $\homeocantor$ is an isometry, then $T$ is an isometry.
 \end{enumerate}
\end{theorem}

\begin{theorem}\label{thm:transitive}
    Let $h \in \groupcantor$ be transitive and  $x \in \cantor$ whose orbit is dense in $\cantor$. Then, there exists $\Phi = (\varphi^L, \varphi^U)$ such that $\fence$ is a Lelek fence and $\hat{h} \in \groupfence$ such that 
    \begin{enumerate}
        %\item[a)] $\varphi (x) > 0 $ for all $x \in orb(h,x)$,
        \item[a)] $h$ is the canonical factor of $\hat{h}$, and
        \item [b)] $(x, \varphi^U(x))$ is a transitive point of $\hat{h}$.
    \end{enumerate}
\end{theorem}

\begin{theorem}\label{thm:mixing}
Let $h:\cantor \rightarrow \cantor$ be a mixing homeomorphism satisfying the following conditions:
for every pair of non-empty open sets \( U, V \subseteq \cantor \), there exists a  compact set \( K_{U,V}\subset \cantor  \) and $m(U,V) \in \N$ such that:
\begin{itemize}
\item $K_{U,V}$ is nowhere dense in $\cantor$ and $h(K_{U,V} )=K_{U,V}$,
    \item for all \( m \geq m(U,V) \), \( h^m(K_{U,V} \cap U) \cap V \neq \emptyset \).
\end{itemize}

Then $h$ admits a lifting $\hat{h}$ to the Lelek fence which is a mixing homeomorphism.
\end{theorem}

\begin{proof}
We let $\varphi^L=0$ and 
    we will define a sequence of partitions $\{G_n\}$  of $\cantor$ and a sequence of functions $\{\varphi^U_n\}$, so that $\fsystem$ will be an  $\cF$-system associated with $h$. In particular, $
    \Psi_n:G_{n+1} \rightarrow G_n$ is the  containment map and $ \oa{uv} \in G_n$ if and only if $h(u) \cap v \neq \emptyset$.  This system will satisfy Condition $\Gamma$ and $\fence$ will be a Lelek fence. Applying Theorem~\ref{thm:mainF}, we will obtain $\hat h$ that is a homeomorphism of $\fence$. Then, we will verify that $\hat{h}$ is mixing.

     Let $G_1= \{\cantor\}$ and $\varphi^U_1:G_1 \rightarrow (0,1]$ be defined as $\varphi^U_1(\cantor) =1$. By hypothesis, let $K_1=K_{\cantor,\cantor}$ and let $m_1 = m(\cantor,\cantor)$.
     
Suppose we are at stage $n$,  compact set $K_n$, $\varphi^U_{n}$, integer $m_n$ and  clopen partitions $G_n$  have been defined so that the following conditions hold.
        \begin{enumerate}[label=(\alph*)]
            \item $h(K_n) = K_n$ and $K_n$ is nowhere dense in $\cantor$.
    \item For all $g, g' \in G_n$ and all $m \ge m_n$ we have that $h^m(K_n \cap g ) \cap g' \neq \emptyset$.
    \item  For all $ g\in G_n$  with $g \cap K_{n-1}  \neq \emptyset$ and $g' \in G_{n-1}$ with $g \subseteq g'$ we have that  $ \varphi^U_{n}(g)= \varphi^U_{n-1}(g')$.
        \end{enumerate}

        We now describe how to construct $K_{n+1}$, $\varphi^U_{n+1}$, $m_{n+1}$ and $G_{n+1}$. 

    
  Let \[  N =  \max  \left \{ \frac{\varphi^U_n(g_1)}{\varphi^U_n(g_2)}: \oa{g_1g_2} \textit { or } \oa{g_2g_1} \in G_n \right \}.\]
    
    Choose $0 < r <1$ so that each of $(1-r)N, |1-r^{-1}|N$ is less than $2^{-(n+1)}$. Now choose $ t \in \N$ sufficiently large so that $r^{t-1}< 2^{-(n+1)}$. Observe that  $\{\varphi^U_n(g), r\varphi^U_n(g), r^2 \varphi^U_n(g), \ldots,\\
    r^{t-1} \varphi^U_n(g)\}$ is $2^{-(n+1)}$ dense in $[0, \varphi^U_n(g)]$, for all $g \in G_n$.

As $h$ is transitive, there is $x \in \cantor$ such that $\orb(h,x)$ is dense in $\cantor$ and  $\{h^j(x): j \in \Z \} \cap K_n = \emptyset$. Let $x_j = h^j(x)$, $j \in \Z$.
Observe that for all $\oa{g_1g_2} \in G_n$, there are infinitely $j$'s  such that $x_j \in g_1$ and $x_{j+1} \in g_2$. This indeed holds since $\{x_j\}_{j \ge 0}$ is dense in $\cantor$. Using this fact, we may choose $\ell\in \N$ and an increasing sequence of integers $\{m_i\}_{i=-t}^t$ with $m_0 =0$,   $m_{-t} = -\ell$, $m_t = \ell$,  so that for all    $\oa{g_1g_2} \in G_n$ and for all $-t \le i < t$  there exists $j$ with $m_i \le j < m_{i+1}$ such that $x_j \in g_1$ and $x_{j+1} \in g_2$. In other words, intuitively speaking, for all $-t \le i < t$, $\{x_{m_i}, \ldots x_{m_{i+1}-1}\}$ realizes all the edges of $G_n$.
     
 Now choose  a clopen set $U$, containing $x_0$ such that the collection $\{h^j(U):-\ell \le j \le \ell \}$ is pairwise disjoint, refines $G_n$ and is disjoint from $K_n$. We may do this as $\{x_n\} \cap K_n = \emptyset$. Let $H_i =\{h^j(U): m_i \le j <m_{i+1} \}$ for $-t \leq i<t$. Finally, enlarge the collection $\cup_{i=-t}^{t-1}H_i$ to form a clopen partition $G_{n+1}$ which refines $G_n$ and such that $mesh (G_{n+1})< 2^{-(n+1)}$. The definition of $G_{n+1}$ is complete.
 
    We now define $\varphi^U_{n+1}: G_{n+1} \rightarrow (0,1]$. Let $g \in G_{n+1}$ with $g \subseteq g'$, $g' \in G_n$. Then,
    \[ \varphi^U_{n+1}(g) = \begin{cases} 
          \varphi^U_n(g') r^{t-1-i }&  \textit { if }  \ \ g \in H_{i},\ \ \ 
           0 \le i < t  \ \ \ \\
           \varphi^U_n(g') r^{t-|i| }&  \textit { if }  \ \ g \in H_{i},  \ \ \ 
           -t \le i < 0 \\
          \varphi^U_n(g')  & \textit{ otherwise }
       \end{cases}
    \]
    In particular, if $g \cap K_n  \neq \emptyset$, then $ \varphi^U_{n+1}(g)= \varphi^U_n(g')$.

  Let  $K_{n+1} = \cup_{g_1,g_2 \in  G_{n+1}} K_{g_1,g_2}$ and $m_{n+1} > \max\{m_n, m(g,g'): g,g' \in G_n \}$. Then, $K_{n+1}$ and $m_{n+1}$ have properties (a)-(c) hold at stage $n+1$ and the induction step $n+1$ is complete. 

We claim that the following conditions hold at stage $n+1$.
    \begin{enumerate}
        %\item $\fsystem$ is an $\cF$-system,
        \item for all $ g\in G_{n+1}$, if $g\cap K_n \neq \emptyset$ then $\varphi^U_{n+1}(g) = \varphi^U_{n}(g') $
        where $g' \in G_n$ is such that $g \subseteq g'$.
        \item if $\oa{g_1g_2}$ or  $\oa{g_2g_1}\in G_{n+1}$, then 
        \[  \left | \frac{\varphi^U_{n+1}(g_1)}{\varphi^U_{n+1}(g_2)}  -\frac{\varphi^U_n(g_1')}{\varphi^U_n(g_2')}\right | < 2^{-(n+1)}
        \]
        where $g_i \subseteq g_i'$, $g_i' \in G_n$ for $i=1,2$,
 
        \item For all $g \in G_n$ and $ 0\le s < t$, there is $g' \in G_{n+1}$ such that $\varphi^U_{n+1} (g')=r^s\varphi^U_n(g)$.
        
        %\item For all $\oa{g_1g_2} \in G_n$ and $ 0\le s < t$ there is $\oa{g_1'g_2'} \in G_{n+1}$ with $g_i'\subseteq g_i$ such that $\varphi^U_{n+1} (g_i')=r^s\varphi^U_n(g_i)$.
    \end{enumerate}

These three conditions are verified analogously as in Theorem~\ref{thm:transitive}.    

That $\fsystem$ is an $\cF$-system follows from the definition of $\varphi^U_{n+1}$. 
Finally, define $\varphi^U = \lim \varphi^U_n$. Condition 2. above implies that the $\Gamma_n < 2^{-(n+1)}$ and hence  Condition $\Gamma$ holds. Furthermore, similarly as in Theorem~\ref{thm:transitive}, Condition 3. implies that $\eta^+_n < 2^{-(n+1)}$. Therefore, by Theorem~\ref{thm:fenceconst} we have that $\fence$ is a Lelek Fence. Thus we have that  $\fsystem$ is an $\cF$-system satisfying Condition $\Gamma$ with $\{(x, \varphi^U(x)\}$ is dense in $\fence$.
Now applying Theorem~\ref{thm:mainF}, we obtain that the resulting $\hat{h}$ is a homeomorphism of $\fence$ whose canonical factor is $h$. 

Now we observe that $\hat h$ is mixing. First, by construction, we have that $\varphi ^U(y) = \varphi^U_n (y)$ for all $y \in K_n$. Recall that the topology on $\fence$ is inherited from the product topology on $\cantor \times [0,1]$. Let $u_i \times (a_i,b_i)$ be open in  $\cantor \times [0,1]$ such that $ u_i \times (a_i,b_i) \cap \fence  \neq \emptyset$, $i=1,2$. It will suffice to show that there exists $n \in \N$, for  all  $m \ge m_n$ we have that 
\[
\hat{h}^{m} (u_1 \times (a_1,b_1) \cap \fence) \cap (u_2 \times (a_2,b_2) \cap \fence)  \neq \emptyset.
\]
To this end, let  $y_i \in \cantor$ such that $(y_i, \varphi^U(y_i)) \in \fence\cap (u_i \times (a_i,b_i))$, $i=1,2$. As $\varphi^U$ is upper semicontinuous, we have that $\varphi^U ((-\infty, b_i))$ is an open subset of $\cantor$ containing $y_i$, $i=1,2$. As $\varphi^U$ is the limit of $\{\varphi^U_n\}$, there exists $n\in \N$ such that  we have $g_i \in G_n$ with $y_i \in g_i \subseteq u_i$ has the property that $g_i \times \{\varphi^U_n(g_i)\} \subseteq u_i \times (a_i,b_i)$ for $i=1,2$. By condition (b), for all  $m \ge m_n$, there is $y\in K_n \cap g_1$ such that  $h^m(y) \in g_2$. Note that 
\[(y, \varphi^U_n(y)) \in u_1 \times (a_1,b_1) 
\]
and 
\[(h^m(y), \varphi^U_n(h^m(y))) \in u_2 \times (a_2,b_2).
\]
As $y \in K_n$, we have that $\varphi^U(y) =\varphi^U_n(y)$ and $h(K_n) =K_n$, we have that $\varphi^U(h^m(y)) =\varphi^U_n(h^m(y))$. Hence, we have 
\[ (y, \varphi^U(y)) = (y, \varphi^U_n(y)) \in u_1 \times (a_1,b_1) \]
and \[ (h^m(y), \varphi^U(h^m(y))) = (h^m(y), \varphi^U_n(h^m(y))) \in u_2 \times (a_2,b_2). \]  Note that \[(y, \varphi^U(y)) \in (u_1 \times (a_1,b_1) \cap \fence) 
\]
and  
\[\hat{h}^m ((y, \varphi^U(y))) = (h^m(y), \varphi^U(h^m(y)))  \in (u_2 \times (a_2,b_2) \cap \fence),
\]
verifying that $\hat{h}$ is mixing and completing the proof.
\end{proof}

\end{document}
